\clearpage
\begin{proofbox}
\textbf{共同准备：弦方程和相切方程，都从基础公式推出来}

设
\[
A_1=P(a),\quad A_2=P(b),\quad A_3=P(c),\quad
P(t)=\left(\frac{t^2}{k},t\right).
\]
三点两两不同，故 $a,b,c$ 也两两不同。

\textbf{第 1 步：写弦，不用斜率。} 对任意 $u\ne v$，把 $P(u),P(v)$ 分别代入
\[
kx-(u+v)y+uv=0
\]
都得到 $0$，所以它就是过两点的直线。此式没有除以 $u+v$，竖直弦也包括在内。

\textbf{第 2 步：直线和圆何时相切？} 圆心 $M(h,0)$ 到上述直线的距离是
\[
d=\frac{|kh+uv|}{\sqrt{k^2+(u+v)^2}}.
\]
因此直线与半径为 $r$ 的圆相切，当且仅当
\begin{equation}
\boxed{(kh+uv)^2=r^2[k^2+(u+v)^2].}\tag{T}
\end{equation}
\begin{center}
% k=r=1, u+v=4/3, uv=-11/3: chord 3x-4y-11=0.
\begin{tikzpicture}[x=1.4cm,y=1.4cm,>=Stealth,font=\small]
\coordinate (M) at (2,0);\coordinate (Q) at ({11/3},0);
\coordinate (R) at (5,1);\coordinate (H) at (2.6,-.8);
\draw[teal!75!black,thick,fill=teal!8] (M) circle(1);
\draw[gray!55,->] (.6,0)--(5.65,0) node[right] {$x$};
\fill[orange!12] (M)--(Q)--(R)--cycle;
\draw[orange!85!black,very thick] (2.2,-1.1)--(5.4,1.3);
\draw[black!65] (M)--(R) (M)--(Q);
\draw[teal!75!black,thick,dashed] (M)--(H) node[midway,left] {$d$};
\draw pic[draw,angle radius=6pt] {right angle=M--H--Q};
\foreach \p in {M,Q,R,H} \fill (\p) circle(1.5pt);
\node[above left] at (M) {$M(h,0)$};
\node[below right] at (Q) {$Q(-uv/k,0)$};
\node[above left] at (R) {$R=Q+(u+v,k)$};
\node[below right] at (H) {$H$};
\node[orange!85!black,align=center] at (4.4,-1.25) {弦所在直线\\$kx-(u+v)y+uv=0$};
\node[align=center] at (3,1.85) {同一三角形：以 $MQ$ 为底，高为 $k$；以 $QR$ 为底，高为 $d$};
\node[align=center] at (3,-2.55) {$|MQ|\,k=|QR|\,d$，故 $d=\dfrac{|kh+uv|}{\sqrt{k^2+(u+v)^2}}$};
\end{tikzpicture}

\end{center}
图中用同一三角形的两种面积求距离：在弦上取 $Q=(-uv/k,0)$，另取 $R=Q+(u+v,k)$；由 $\triangle MQR$ 的两种“底乘高除以二”的算法，分别以 $MQ$ 和 $QR$ 为底，就得到上面的 $d$。

\end{proofbox}
\clearpage
\begin{proofbox}
\textbf{共同准备（续）：把两条相切条件写成代数式}

\textbf{第 3 步：用圆心条件整理成同一种二次式。} 因为 $kh=kr+r^2$，把 $u=a,v=b$ 代入 (T) 并展开，得到
\begin{equation}
(a^2-r^2)b^2+2kra b+r^2(2kr+r^2-a^2)=0.\tag{1}
\end{equation}
把 $v$ 改成 $c$，同理
\begin{equation}
(a^2-r^2)c^2+2kra c+r^2(2kr+r^2-a^2)=0.\tag{2}
\end{equation}
以下四种方法都只用到这两个式子。\textbf{注意：}硬算方法不要求预先认出“同一个方程的两根”。
\end{proofbox}

\clearpage
\begin{proofbox}
\textbf{方法一：最低门槛——你熟悉的硬算，相减求和、相加求积}

\begin{center}
% Exact non-symmetric example a=1/2, b,c=(2+-sqrt37)/3.
\begin{tikzpicture}[x=.9cm,y=.9cm,>=Stealth,font=\small]
\pgfmathsetmacro{\bb}{(2+sqrt(37))/3}
\pgfmathsetmacro{\cc}{(2-sqrt(37))/3}
\coordinate (A1) at (.25,.5);\coordinate (A2) at ({\bb*\bb},\bb);
\coordinate (A3) at ({\cc*\cc},\cc);\coordinate (M) at (2,0);\coordinate (H) at (2.6,-.8);
\draw[gray!55,->] (-.5,0)--(8.1,0) node[right] {$x$};
\draw[gray!55,->] (0,-1.8)--(0,3.2) node[above] {$y$};
\draw[blue!70!black,thick,domain=-1.65:2.78,samples=110,variable=\t] plot ({\t*\t},\t);
\draw[teal!75!black,thick,fill=teal!9] (M) circle(1);
\draw[orange!85!black,thick] (A2)--(A1)--(A3);
\draw[orange!85!black,very thick] (A2)--(A3);
\draw[teal!75!black,dashed] (M)--(H) node[midway,above right] {$r$};
\draw pic[draw,angle radius=5pt] {right angle=M--H--A2};
\foreach \p in {A1,A2,A3} \fill[orange!85!black] (\p) circle(1.8pt);
\fill[teal!75!black] (M) circle(1.5pt) node[above] {$M$};
\node[above left] at (A1) {$A_1=P(a)$};
\node[above left] at (A2) {$A_2=P(b)$};
\node[below left] at (A3) {$A_3=P(c)$};
\node[anchor=west,align=left] at (5,-.7) {第三条弦只需要\\[4pt]和 $b+c$、积 $bc$};
\node[align=center] at (4,-3.05) {$a=\frac12,\ b+c=\frac43,\ bc=-\frac{11}{3}$\\[7pt]$A_2A_3:\ x-\frac43y-\frac{11}{3}=0$};
\end{tikzpicture}

\end{center}
图中第三条弦由和、积确定；下面只求这两个量，不逐个解点。

\textbf{第 1 步：两式相减。} (1) 减 (2)，相同常数自动消掉：
\[
(a^2-r^2)(b^2-c^2)+2kra(b-c)=0.
\]
用平方差公式 $b^2-c^2=(b-c)(b+c)$，且 $b\ne c$，可以约去 $b-c$：
\begin{equation}
\boxed{(a^2-r^2)(b+c)=-2kra.}\tag{3}
\end{equation}
由此还能确认 $a^2\ne r^2$：否则左边等于 $0$，右边 $-2kra\ne0$（$k,r>0$，$a=\pm r\ne0$），矛盾。

\end{proofbox}
\clearpage
\begin{proofbox}
\textbf{方法一（续）：已有和，再求积}

\textbf{第 2 步：两式相加。} (1) 加 (2)：
\[
\begin{aligned}
&(a^2-r^2)(b^2+c^2)+2kra(b+c)\\
&\qquad+2r^2(2kr+r^2-a^2)=0.
\end{aligned}
\]
我们\textbf{已经有和，目标缺积}，所以写
\[
\boxed{b^2+c^2=(b+c)^2-2bc}.
\]
代入后，$ (a^2-r^2)(b+c)^2 $ 由 (3) 正好等于 $-2kra(b+c)$，与原有的 $+2kra(b+c)$ 抵消，剩下
\[
-2(a^2-r^2)bc+2r^2(2kr+r^2-a^2)=0.
\]
所以
\begin{equation}
\boxed{(a^2-r^2)bc=r^2(2kr+r^2-a^2).}\tag{4}
\end{equation}
这一步是本题硬算的\textbf{关键突破}：\textbf{相减得到和，相加利用平方和得到积}，不用解出 $b,c$。
\end{proofbox}

\clearpage
\begin{proofbox}
\textbf{方法一（续）：把“和、积”放回第三条弦}

因为 $A_2A_3$ 是 $kx-(b+c)y+bc=0$，只需按 (T) 检查。

由 (4) 和 $kh=kr+r^2$，
\[
\begin{aligned}
kh+bc&=kr+r^2+\frac{r^2(2kr+r^2-a^2)}{a^2-r^2}\\
&=\frac{kr(a^2+r^2)}{a^2-r^2}.
\end{aligned}
\]
由 (3) 算出
\[
\begin{aligned}
k^2+(b+c)^2
&=k^2+\frac{4k^2r^2a^2}{(a^2-r^2)^2}\\
&=\frac{k^2\bigl[(a^2-r^2)^2+4r^2a^2\bigr]}{(a^2-r^2)^2}\\
&=\frac{k^2(a^2+r^2)^2}{(a^2-r^2)^2}.
\end{aligned}
\]
于是
\[
\boxed{(kh+bc)^2=r^2[k^2+(b+c)^2]}
\]
恰好符合 (T)，所以\textbf{$A_2A_3$ 与 $\Gamma$ 相切}。

\medskip
\textbf{回看方法一配图：}求出的和、积确定第三条弦；上式说明圆心到它的距离正好是 $r$，完成闭合。
\end{proofbox}

\clearpage
\begin{proofbox}
\textbf{\textcolor{CProof}{方法二：最快——把两个相切条件看成同一个方程的两根}}

\begin{center}
% f(t)=-3t^2/4+t+11/4; roots c,b and axis (b+c)/2=2/3.
\begin{tikzpicture}[x=1.45cm,y=.75cm,>=Stealth,font=\small]
\pgfmathsetmacro{\bb}{(2+sqrt(37))/3}
\pgfmathsetmacro{\cc}{(2-sqrt(37))/3}
\draw[gray!55,->] (-2,0)--(3.3,0) node[right] {$t$};
\draw[gray!55,->] (0,-1.7)--(0,4.1);
\draw[blue!70!black,very thick,domain=-1.75:3.08,samples=100,variable=\t] plot (\t,{-.75*\t*\t+\t+2.75});
\draw[teal!75!black,dashed] ({2/3},-.55)--({2/3},{37/12});
\fill[orange!85!black] (\bb,0) circle(2pt) node[above right=4pt] {$b$};
\fill[orange!85!black] (\cc,0) circle(2pt) node[above left=4pt] {$c$};
\node[below,align=center] at ({2/3},-.6) {对称轴 $t=\dfrac{b+c}{2}$};
\node[align=center] at (.6,5) {固定 $a=\frac12$，取 $k=r=1$\\[4pt]$f(t)=-\frac34t^2+t+\frac{11}{4}$};
\node[align=center] at (.6,-3.4) {两个横轴交点 $\longleftrightarrow$ 两个不同参数 $b,c$\\[7pt]韦达定理直接给出 $b+c=\frac43,\ bc=-\frac{11}{3}$};
\end{tikzpicture}

\end{center}
图把两条切圆弦对应的参数画成同一条函数图像的两个零点；不必分别求根。

\textbf{识别信号：}固定 $A_1=P(a)$，不同的 $A_2,A_3$ 满足\textbf{同一条相切条件}。将 (1)、(2) 写成同一个未知数 $t$ 的方程：
\[
(a^2-r^2)t^2+2krat+r^2(2kr+r^2-a^2)=0.
\]
这个方程的两个不同的根是 $b,c$。如果 $a^2=r^2$，二次项消失而一次项 $2kra\ne0$，方程最多只有一个根，与 $b\ne c$ 矛盾。因此 $a^2-r^2\ne0$，可以用韦达定理一次得到
\[
\boxed{b+c=-\frac{2kra}{a^2-r^2}},\qquad
\boxed{bc=\frac{r^2(2kr+r^2-a^2)}{a^2-r^2}}.
\]
再算
\[
kh+bc=\frac{kr(a^2+r^2)}{a^2-r^2},\quad
k^2+(b+c)^2=\frac{k^2(a^2+r^2)^2}{(a^2-r^2)^2}.
\]
故满足 (T)，$A_2A_3$ 与圆相切。

\textbf{回看本方法配图：}两个零点对应 $b,c$；韦达定理把两点的信息压缩成第三条弦所需的和与积。
\end{proofbox}

\clearpage
\begin{proofbox}
\textbf{\textcolor{CProof}{方法三：步骤最少——相减后回代，不再相加}}

\begin{center}
% Positive representative b>0,c<0,b+c>0. This is an area model,
% not a sign assumption in the general algebraic proof.
\begin{tikzpicture}[x=1.15cm,y=1.15cm,>=Stealth,font=\small]
\pgfmathsetmacro{\bb}{(2+sqrt(37))/3}
\draw[fill=orange!18,draw=none] (0,0) rectangle ({4/3},\bb);
\draw[fill=teal!18,draw=none] ({4/3},0) rectangle (\bb,\bb);
\draw[thick] (0,0) rectangle (\bb,\bb);
\draw[thick,dashed] ({4/3},0)--({4/3},\bb);
\draw[<->] (-.4,0)--(-.4,\bb) node[midway,left] {$b$};
\draw[<->] (0,-.4)--({4/3},-.4) node[midway,below] {$b+c$};
\draw[<->] ({4/3},-.4)--(\bb,-.4) node[midway,below] {$-c$};
\node[align=center] at ({2/3},{\bb/2}) {减去\\$b(b+c)$};
\node[align=center] at ({(\bb+4/3)/2},{\bb/2}) {剩下\\$-bc$};
\node[align=left,anchor=west] at (3.25,1.6) {原正方形面积 $b^2$\\[7pt]减去左侧矩形\\[7pt]即剩右侧矩形};
\node[align=center] at (2.7,-1.8) {$b^2-b(b+c)=b[b-(b+c)]=-bc$\\[7pt]图取 $b>0,\ c<0,\ b+c>0$；恒等式对任意实数均成立};
\end{tikzpicture}

\end{center}
先看图中“正方形减矩形”的剩余面积，再把同一个恒等式用于一般参数。

由 (1) 减 (2) 已得
\[
(a^2-r^2)(b+c)=-2kra.
\]
将 $2kra=-(a^2-r^2)(b+c)$ 乘以 $b$，代回 (1)：
\[
\begin{aligned}
&(a^2-r^2)b^2-(a^2-r^2)b(b+c)\\
&\qquad+r^2(2kr+r^2-a^2)=0.
\end{aligned}
\]
利用最基础的恒等式
\[
\boxed{b^2-b(b+c)=-bc},
\]
立即得到
\[
(a^2-r^2)bc=r^2(2kr+r^2-a^2).
\]
和、积齐全，代入 (T) 即可（完整的最后验算见方法一续）。

\textbf{作用：}少写一条相加方程；但它并不一定比“相减再相加”更容易在考场想到。\textbf{两者都属于纯代数硬算。}
\end{proofbox}

\clearpage
\begin{proofbox}
\textbf{\textcolor{CProof}{方法四：常规——直接联立第三条直线和圆，判别式等于零}}

\begin{center}
% All three circles have center (2,0), radius 1, equal coordinate scales.
\begin{tikzpicture}[x=1cm,y=1cm,>=Stealth,font=\small]
\foreach \shift/\line/\title/\disc in {0/2.6/两个交点/\Delta>0,4/3/一个交点/\Delta=0,8/3.4/没有交点/\Delta<0}{
\begin{scope}[xshift=\shift cm]
\draw[teal!75!black,thick,fill=teal!9] (2,0) circle(1);
\draw[orange!85!black,very thick] (\line,-1.5)--(\line,1.5);
\fill[teal!75!black] (2,0) circle(1.6pt) node[left] {$M$};
\node at (2,2) {\title};
\node at (2,-2) {$\disc$};
\end{scope}}
\fill[orange!85!black] (2.6,.8) circle(2pt) (2.6,-.8) circle(2pt) (7,0) circle(2pt);
\node[right] at (7,0) {$H$};
\node[align=center] at (6,-2.95) {中图是 $a=0$ 的闭合例：第三条弦 $x=3$\\[5pt]代入 $(x-2)^2+y^2=1$ 得 $y^2=0$，唯一交点 $H=(3,0)$};
\end{tikzpicture}

\end{center}
图中固定同一个圆，只平移直线；本方法要证明交点数恰好是中间的一。

由 (3)、(4) 记
\[
S=b+c=-\frac{2kra}{a^2-r^2},\qquad
T=bc=\frac{r^2(2kr+r^2-a^2)}{a^2-r^2}.
\]
第三条弦为 $kx-Sy+T=0$，故 $x=(Sy-T)/k$。将其代入圆 $(x-h)^2+y^2=r^2$，两边同乘 $k^2$：
\[
(Sy-T-kh)^2+k^2y^2=k^2r^2.
\]
展开后得到
\[
(S^2+k^2)y^2-2S(T+kh)y+(T+kh)^2-k^2r^2=0.
\]
此方程的判别式为
\[
\begin{aligned}
\Delta&=4S^2(T+kh)^2\\
&\quad-4(S^2+k^2)\bigl[(T+kh)^2-k^2r^2\bigr]\\
&=4k^2\bigl[r^2(S^2+k^2)-(T+kh)^2\bigr].
\end{aligned}
\]
方法一末尾已经算出 $(T+kh)^2=r^2(S^2+k^2)$，所以 $\Delta=0$。因 $S^2+k^2>0$，交点只有一个，第三条弦确实与圆相切。

\end{proofbox}
\clearpage
\begin{proofbox}
\textbf{补充另一圆心位置（仍是整条直线）：}若 $h=r^2/k-r$，相切式 (T) 对应
\[
(a^2-r^2)t^2-2krat+r^2(r^2-2kr-a^2)=0.
\]
其两个不同的解 $b,c$ 满足
\[
 b+c=\frac{2kra}{a^2-r^2},\quad
 bc=\frac{r^2(r^2-2kr-a^2)}{a^2-r^2}.
\]
直接计算得
\[
kh+bc=-\frac{kr(a^2+r^2)}{a^2-r^2},\quad
k^2+(b+c)^2=\frac{k^2(a^2+r^2)^2}{(a^2-r^2)^2};
\]
仍满足 (T)。注意该分支中切点可能在\textbf{线段延长线}上，不能直接说是三角形内切圆。
\end{proofbox}
